\section{Training configurations}
\label{app:training}

This appendix documents the architecture and training configuration for
each of the four architectures (MAMNet, OGLANet, DINOv3, CrossEarth)
and the eight adaptation methods (FDA, SegDeSiC, IIM, ISW, MRFP+, FADA,
BSM, SAR).  SegDeSiC, IIM, ISW, MRFP+ and FADA are not run on
CrossEarth, whose query-based Mask2Former decoder has no per-pixel head
these methods are defined on.  Each adaptation block lists only those
hyperparameters that differ from the corresponding base configuration;
values matching the base are omitted.  

\subsection{MAMNet}
\label{sec:mamnet-base}

\textbf{Architecture.}
ResNet-34~\cite{he2016deep} backbone with pretrained weights
\texttt{ResNet34\_Weights.IMAGENET1K\_V1} (torchvision); \texttt{conv1}
stride changed from 2 to 1 and \texttt{maxpool} replaced by
\texttt{nn.Identity}, so the encoder returns five feature maps at
channels $\{64,128,256,512,512\}$ with output strides
$\{1,2,4,8,16\}$ relative to the $384\times384$ input
(\texttt{feat5}: $512{\times}24{\times}24$).
The MSCAF module operates on \texttt{feat5} before the decoder.
The decoder comprises four \texttt{DecoderBlock} stages
($512{\to}256{\to}128{\to}64{\to}64$); each block applies Criss-Cross
Attention (CCA) to the incoming skip-connection feature, bilinearly
upsamples the previous decoder map, concatenates, and refines with
two $3{\times}3$-BN-ReLU convolutions.  CCA uses $1{\times}1$
projections to $C/8$ channels for $Q,K$ and $C$ for $V$; row- and
column-wise attention maps are summed and gated by a learnable scalar
$\gamma$ (init.\ 0): $\text{out} = \gamma(H{+}W) + x$.  Three
auxiliary branches (AUX1--3) attach to decoder outputs at $256/128/64$
channels (Conv$3{\times}3$-BN-ReLU-Dropout(0.3)-Conv$1{\times}1$,
bilinearly upsampled to input resolution).
Input is $384{\times}384$ four-channel RGBC; the contrast-channel
\texttt{conv1} weight is initialised as the channel-mean of the
pretrained RGB weights.

\begin{table}[h]
\centering
\caption{MAMNet training configuration.
$^\dagger$Shell script overrides the Python-file default of
$1{\times}10^{-3}$.}
\label{tab:mamnet-config}
\small
\begin{tabular}{@{}ll@{}}
\toprule
\textbf{Hyperparameter} & \textbf{Value} \\
\midrule
\multicolumn{2}{@{}l}{\textit{Optimisation}} \\
Optimizer        & Adam \\
Learning rate    & $1{\times}10^{-4}$$^\dagger$ \\
Weight decay     & $1{\times}10^{-4}$ \\
Batch size       & 8 \\
Epochs           & 100 \\
LR scheduler     & \texttt{ReduceLROnPlateau}
                   (mode=\texttt{max}, factor=0.5, patience=3) \\
Decision metric  & Per-image tolerant mIoU ($\pm$2\,px) \\
Early stopping   & patience = 10 epochs \\
\midrule
\multicolumn{2}{@{}l}{\textit{Loss}} \\
Main loss        & \texttt{nn.CrossEntropyLoss} (2-class, unweighted) \\
Total loss       & $\mathcal{L}_\text{main} + 0.4\cdot
                    \tfrac{1}{3}\sum_{k=1}^{3}\mathcal{L}_{\text{aux}_k}$ \\
\midrule
\multicolumn{2}{@{}l}{\textit{Data}} \\
Input size       & $384{\times}384$ \\
Channels         & 4 (RGBC) \\
Train aug.       & H-flip $p{=}0.5$, V-flip $p{=}0.5$ \\
Val/test aug.    & None (resize + ImageNet normalisation only) \\
Post-processing  & Remove connected components $<10$\,px \\
\midrule
\multicolumn{2}{@{}l}{\textit{Infrastructure}} \\
Hardware         & 1$\times$ NVIDIA A100 (NCSA Delta, \texttt{gpuA100x4}) \\
\bottomrule
\end{tabular}
\end{table}

\subsection{MAMNet + FDA}

\textbf{MAMNet\,+\,FDA.}
FDA~\cite{fda} is applied at the input stage, before resize and
normalisation: each source training image (uint8 PIL $\to$ numpy) is
passed through \texttt{FDATransform} on CPU.  Per-channel 2-D FFT
separates amplitude and phase; the source amplitude rectangle of
half-width $b{=}\lfloor\min(H,W)\cdot L\rfloor$ centred on the shifted
spectrum is replaced element-wise by the corresponding target
rectangle (no blending), and the image is reconstructed from the
mutated amplitude with the source phase
$\hat{\mathbf{s}} = \operatorname{Re}\bigl(\mathcal{F}^{-1}(\hat{\mathcal{A}}_s
e^{i\mathcal{P}_s})\bigr)$, clipped to $[0,255]$.
With $L{=}0.01$ and $\min(H,W){=}384$, the swapped region is $6{\times}6$
per channel.  Validation and test images are unmodified.  Run in LOCO
mode only; the target pool is up to 100 unlabelled images sampled
uniformly per iteration from the held-out city's \texttt{train/images}.

\begin{table}[h]
\centering
\caption{MAMNet\,+\,FDA hyperparameters; settings matching
Table~\ref{tab:mamnet-config} are omitted.}
\label{tab:mamnet-fda}
\small
\begin{tabular}{@{}ll@{}}
\toprule
\textbf{Hyperparameter} & \textbf{Value} \\
\midrule
FDA low-freq.\ ratio $L$ & 0.01 \\
Target pool size         & $\leq 100$ images \\
Target domain            & Held-out city \texttt{train/images} (unlabelled) \\
Application point        & Input (pre-resize, train only) \\
Target sampling          & Uniform random per iteration \\
Early-stopping patience  & 15 epochs \\
Modes run                & LOCO only \\
\bottomrule
\end{tabular}
\end{table}

\subsection{MAMNet + SegDeSiC}

\textbf{MAMNet\,+\,SegDeSiC.}
SegDeSiC~\cite{segdesicnet} attaches a geographic coordinate
branch to \texttt{feat5} ($512{\times}24{\times}24$): global average
pool $\to$ MLP
($512{\to}256{\to}256{\to}40$, ReLU + Dropout(0.3) between layers)
$\to$ $\ell_2$ normalisation, predicting a 40-d GRID encoding of
$(\varphi,\lambda)$ produced by sphere projection
($x{=}\cos\varphi\cos\lambda$, $y{=}\cos\varphi\sin\lambda$) and
$[\sin(\omega_k x),\cos(\omega_k x),\sin(\omega_k y),\cos(\omega_k y)]$
across 10 log-spaced wavelengths
$\lambda_k\in[10^{-4},10^{4}]$ ($\omega_k{=}2\pi/\lambda_k$).
The domain loss is
$\mathcal{L}_{\text{dom}}{=}1-\cos\langle\hat{e},e^{*}\rangle$,
batch-mean.
In LOCO mode the loss is computed on both the labelled source batch
and an unlabelled target batch (held-out city, no segmentation labels);
both terms are summed before the $\alpha$ scaling.  No coordinates are
passed at validation or test time.  

\begin{table}[h]
\centering
\caption{MAMNet\,+\,SegDeSiC hyperparameters; settings matching
Table~\ref{tab:mamnet-config} are omitted.
$^\dagger$No submit script provided; \texttt{argparse} default may have
been overridden at launch.}
\label{tab:mamnet-segdesic}
\small
\begin{tabular}{@{}ll@{}}
\toprule
\textbf{Hyperparameter} & \textbf{Value} \\
\midrule
\multicolumn{2}{@{}l}{\textit{Geographic branch}} \\
Attachment point         & \texttt{feat5} ($512{\times}24{\times}24$) \\
GRID scales / range      & 10 / $[10^{-4},10^{4}]$ log-spaced \\
Encoding dim             & 40 ($4{\times}10$) \\
MLP hidden dim / dropout & 256 / 0.3 (both layers) \\
\midrule
\multicolumn{2}{@{}l}{\textit{Loss}} \\
Domain loss              & $1-\cos\langle\hat{e},e^{*}\rangle$, batch mean \\
$\alpha$ (domain weight) & 0.5 \\
Total loss               & $\mathcal{L}_{\text{seg}}
                            + \alpha(\mathcal{L}_{\text{dom,src}}
                                    + \mathcal{L}_{\text{dom,tgt}})$ \\
CE ignore index          & 255 \\
\midrule
\multicolumn{2}{@{}l}{\textit{Optimisation}} \\
Learning rate            & $1{\times}10^{-3}$ \\
Early-stopping patience  & 15 epochs \\
\midrule
\multicolumn{2}{@{}l}{\textit{Data}} \\
Target domain (LOCO)     & Held-out city \texttt{train/images} (unlabelled) \\
Geo.\ metadata           & Per-image lat/lon in \texttt{mapping.json} \\
Modes run                & LOCO only \\
\bottomrule
\end{tabular}
\end{table}

\subsection{MAMNet + IIM}

\textbf{MAMNet\,+\,IIM.}
IIM~\cite{yola} is an input-level preprocessor placed before the
backbone: input RGB is denormalised, log-transformed, and convolved by
$K{=}8$ learnable $5{\times}5$ kernels in three cross-channel pairs
($W{*}\log R-W{*}\log B$, $W{*}\log R-W{*}\log G$,
$W{*}\log G-W{*}\log B$), producing 24 channels.  A FuseConv block
fuses these with the original RGB through two parallel
$3{\times}3$ paths (each $\to$ BN-ReLU), concatenation, and two
$3{\times}3$ convolutions, returning a 3-channel output that replaces
the RGB input to MAMNet.  After every \texttt{optimizer.step()} each
kernel is projected to zero-mean
($W{\leftarrow}W-\overline{W}$) to enforce illumination cancellation
in the log-domain.
A consistency loss
$\mathcal{L}_{\text{II}} = \operatorname{SmoothL1}\!\bigl(
\text{IIM}(I),\,\text{IIM}(\sigma_\gamma I)\bigr)$
with random $\gamma\sim U[0.5,2.0]$ regularises kernel invariance; an
adaptive weight
$\eta = r\cdot\mathcal{L}_{\text{task}}/(\mathcal{L}_{\text{II}}{+}\epsilon)$
holds the II contribution at $r{=}1\%$ of the task loss
(gradients flow only through $\mathcal{L}_{\text{II}}$).
IIM adds 0.008\,M parameters.

\begin{table}[h]
\centering
\caption{MAMNet\,+\,IIM hyperparameters; settings matching
Table~\ref{tab:mamnet-config} are omitted.}
\label{tab:mamnet-iim}
\small
\begin{tabular}{@{}ll@{}}
\toprule
\textbf{Hyperparameter} & \textbf{Value} \\
\midrule
\multicolumn{2}{@{}l}{\textit{IIM module}} \\
Num.\ kernels $K$        & 8 \\
Kernel size              & $5{\times}5$ \\
Cross-channel pairs      & (R,B), (R,G), (G,B) $\to$ 24 features \\
Zero-mean projection     & After every optimizer step \\
\midrule
\multicolumn{2}{@{}l}{\textit{Loss}} \\
II loss                  & SmoothL1 ($\beta{=}1.0$) on IIM features \\
II weighting             & Adaptive, target ratio $r{=}0.01$ \\
Gamma range              & $\gamma\sim U[0.5,2.0]$ \\
Total loss               & $\mathcal{L}_{\text{task}}
                            + \eta\,\mathcal{L}_{\text{II}}$ \\
\bottomrule
\end{tabular}
\end{table}

\subsection{MAMNet + ISW}

\textbf{MAMNet\,+\,ISW.}
ISW~\cite{isw} is a training-only feature-covariance
regulariser; the inference architecture is unchanged.
\emph{Phase 1 (offline)}: each training image and its photometric
augmentation (\texttt{ColorJitter(0.4,0.4,0.4,0.1)}\,+\,\texttt{Gaussian\-Blur}
$p{=}0.5$) are passed through the encoder; the per-instance covariance
of standardised features at \texttt{feat1}, \texttt{feat2},
\texttt{feat3} is computed, and the variance of each covariance
element across the original/augmented pair is averaged over the
training set.  Upper-triangular variance values are partitioned by
$k$-means ($k{=}3$); the top $k{-}m{=}2$ clusters are flagged as
style-sensitive, producing a binary mask $\widetilde M$ per layer.
\emph{Phase 2 (training)}: at each step, forward hooks capture
features at the same three layers; per-image standardised covariance
is masked element-wise with $\widetilde M$ and reduced to
$\mathcal{L}_{\text{ISW}}^{\ell}{=}\|\widetilde M\odot\Sigma_s\|_1$
(per-image mean), then averaged across the three layers.
Total loss is $\mathcal{L}_{\text{seg}}+\lambda\,\mathcal{L}_{\text{ISW}}$.

\begin{table}[h]
\centering
\caption{MAMNet\,+\,ISW hyperparameters; settings matching
Table~\ref{tab:mamnet-config} are omitted.}
\label{tab:mamnet-isw}
\small
\begin{tabular}{@{}ll@{}}
\toprule
\textbf{Hyperparameter} & \textbf{Value} \\
\midrule
\multicolumn{2}{@{}l}{\textit{ISW configuration}} \\
$\lambda$ (ISW weight)         & 0.6 \\
Hooked layers                  & \texttt{feat1, feat2, feat3} \\
$k$-means $k$ / $m$            & 3 / 1 (bottom 1 cluster $=$ content) \\
Photometric aug.\ (mask precomp.)\!\!
  & \texttt{ColorJitter(0.4,0.4,0.4,0.1)} \\
  & + \texttt{GaussianBlur} $p{=}0.5$ \\
Loss reduction                 & Per-image mean, then layer mean \\
Total loss                     & $\mathcal{L}_{\text{seg}}
                                  + \lambda\,\mathcal{L}_{\text{ISW}}$ \\
\bottomrule
\end{tabular}
\end{table}

\subsection{MAMNet + MRFP+}

\textbf{MAMNet\,+\,MRFP+.}
MRFP+~\cite{mrfp} adds three training-only stochastic
perturbations.  HRFP attaches after \texttt{conv1}-\texttt{bn1}-\texttt{relu}
and before \texttt{layer1} (64\,ch, $H/2$): a 4-encoder/4-decoder
overcomplete autoencoder with $\sim$1.2$\times$ bilinear upsampling per
encoder layer; its conv weights are He-reinitialised and BN
$\gamma,\beta$ resampled from $\mathcal{N}(0,0.5^2)$ \emph{at every
forward pass}, and its output is added to the encoder stream under
\texttt{no\_grad}+\texttt{detach}.  NP+ attaches at \texttt{feat4}
($512{\times}48{\times}48$): channel-mean perturbation
$y = \alpha x + \delta(\beta-\alpha)\mu_c$ with $\alpha,\beta\sim
\mathcal{N}(0,1)$ and $\delta$ a normalised cross-batch variance.
HRFP+ adds the max-resolution intermediate feature from the HRFP
encoder to decoder stage \texttt{dec3}; three parameter-free
\texttt{InstanceNorm2d} layers (affine=False) are applied on
\texttt{feat1}/\texttt{feat2}/\texttt{feat3}.  Each module is toggled
independently with $p{=}0.5$.  At inference all modules are disabled
and the model is identical to base MAMNet (HRFP weights add $\sim$74\,K
frozen parameters never saved).

\begin{table}[h]
\centering
\caption{MAMNet\,+\,MRFP+ hyperparameters; settings matching
Table~\ref{tab:mamnet-config} are omitted.}
\label{tab:mamnet-mrfp}
\small
\begin{tabular}{@{}ll@{}}
\toprule
\textbf{Hyperparameter} & \textbf{Value} \\
\midrule
\multicolumn{2}{@{}l}{\textit{Perturbation modules (training only)}} \\
HRFP attach point        & After stem (64\,ch, $H/2$) \\
NP+ attach point         & \texttt{feat4} (512\,ch, $H/32$) \\
HRFP+ attach point       & Decoder \texttt{dec3} (64\,ch) \\
Inst.\ Norm layers       & \texttt{feat1}, \texttt{feat2}, \texttt{feat3}
                           (affine=False) \\
HRFP / NP+ / HRFP+ prob.\ & 0.5 / 0.5 / 0.5 (independent) \\
HRFP arch.               & 4 enc + 4 dec, $\sim$1.2$\times$ scale per layer \\
HRFP BN std $\sigma$     & 0.5 \\
HRFP weights             & Frozen, re-randomised every forward pass \\
\bottomrule
\end{tabular}
\end{table}

\subsection{MAMNet + FADA}

\textbf{MAMNet\,+\,FADA.}
FADA~\cite{fada} freezes the ResNet-34 encoder
(\texttt{requires\_grad=False}; BN in \texttt{eval} mode using
pretrained running statistics) and inserts a frequency adapter at
\texttt{feat3}, \texttt{feat4}, and \texttt{feat5}.  Each FADA block
applies a 2-D Haar DWT $\to$ $\{$LL, LH, HL, HH$\}$;
the LL subband is processed by a low-frequency branch (LoRA tokens
$T_L\in\mathbb{R}^{m\times c}$ with $T{=}AB$, rank $r$; scaled-dot
attention; MLP projection; skip fusion) for content stabilisation;
the concatenated LH/HL/HH are processed by a high-frequency branch
(separate LoRA tokens $T_H\in\mathbb{R}^{3m\times c}$; scaled-dot
attention; \emph{Instance Normalisation applied to the similarity
map}; MLP projection; skip fusion) for style suppression; the four
adapted subbands are recombined by inverse Haar DWT.  Stages 1 and 2
are excluded for memory reasons (spatial dims $384^2$ and $192^2$
yield prohibitive token attention).
Trainable parameters: $\approx$8.4\,M (FADA $\approx$2.5\,M, MSCAF +
decoder + AUX $\approx$5.9\,M); frozen $\approx$21.3\,M.

\begin{table}[h]
\centering
\caption{MAMNet\,+\,FADA hyperparameters; settings matching
Table~\ref{tab:mamnet-config} are omitted.}
\label{tab:mamnet-fada}
\small
\begin{tabular}{@{}ll@{}}
\toprule
\textbf{Hyperparameter} & \textbf{Value} \\
\midrule
\multicolumn{2}{@{}l}{\textit{FADA configuration}} \\
Adapter stages           & \texttt{feat3, feat4, feat5} \\
LoRA rank $r$            & 16 \\
Token length $m$         & 100 \\
LF branch                & LoRA tokens on LL (content) \\
HF branch                & LoRA tokens + IN on similarity map (LH/HL/HH) \\
Decomposition            & Haar DWT (orthogonal, invertible) \\
\midrule
\multicolumn{2}{@{}l}{\textit{Optimisation overrides}} \\
Backbone state           & Frozen (BN in eval mode) \\
Trainable groups         & FADA blocks + MSCAF + decoder + AUX \\
\bottomrule
\end{tabular}
\end{table}

\subsection{OGLANet}
\label{sec:oglanet-base}

\textbf{Architecture.}
OGLANet~\cite{oglanet} is a four-stage encoder--decoder with
dense skip connections.  The encoder consists of five Global-Local
Aware Module (GLAM) stages on top of a ResNet-34 backbone (pretrained
on ImageNet via \texttt{torchvision.models.resnet34}); the original
paper specifies ResNet-101, but ResNet-34 is substituted to match
MAMNet's parameter budget.  Each GLAM wraps one ResNet layer (LFEM)
between two GFEM self-attention modules ($1{\times}1$ Q/K/V
projections, sigmoid-gated, followed by $3{\times}3$ conv with
stride 2 for downsampling or stride 1 for refinement).  GLAM5 is an
additional stage (not in the original figure) that expands 512$\to$1024
channels with a stride-2 $3{\times}3$ conv plus a non-downsampling
GFEM.  The DFFM forms four decoder maps $S^D_1$--$S^D_4$ by
concatenating multi-scale encoder features iteratively downsampled
with stride-2 $3{\times}3$ convs, each followed by a $1{\times}1$
channel-reduction conv (BN-ReLU).  The decoder bilinearly upsamples
DFFM features through Conv$3{\times}3$-BN-ReLU blocks, producing
outputs at $\{12{\times}12, 48{\times}48, 96{\times}96,
384{\times}384\}$.  The OAM produces six $1{\times}1$-conv predictions:
$P_1$--$P_4$ from the four decoder scales, $P_5$ from
$\operatorname{concat}(P_1{:}P_4)$, $P_6$ from
$\operatorname{concat}(P_1{:}P_5)$.  All six are deeply supervised at
training; only $P_6$ is used at inference.  Input is $384{\times}384
{\times}4$ (RGBC); the contrast-channel \texttt{conv1} weight is
initialised as the channel-mean of the pretrained RGB weights.
Output stride 1.  

\textbf{DFFM reimplementation defect.}\label{app:oglanet-dffm}
We reimplemented the DFFM from the published architecture figure, and the
reimplementation makes OGLANet inference depend on the random state and the
batch size.  Its $20$ downsampling convolutions and their batch-normalization
layers are constructed inside \texttt{forward()} on every call and never
registered as modules.  They are absent from the checkpoint, receive fresh
random weights on every forward pass, and stay in training mode under
\texttt{model.eval()}, so they normalize with the statistics of the current
batch.  With a fixed seed, inference is bit-identical across jobs and nodes.
Over repeated unseeded passes of one checkpoint on the same $150$ test images,
the standard deviation of the cell mean reaches $0.42$ mIoU.  Changing the
inference seed among three values moves the cell mean by up to $0.66$ mIoU on
an OGLANet-\sib{} checkpoint and by up to $1.37$ mIoU on a published Vanilla
checkpoint, both on Chicago at high resolution.  Every OGLANet value in the
paper is one inference pass at its run's inference batch size, and its
confidence interval does not include this variation.  We left the module
unchanged because every reported OGLANet number was produced with it, and the
released code contains the same module.

\begin{table}[h]
\centering
\caption{OGLANet training configuration.}
\label{tab:oglanet-config}
\small
\begin{tabular}{@{}ll@{}}
\toprule
\textbf{Hyperparameter} & \textbf{Value} \\
\midrule
\multicolumn{2}{@{}l}{\textit{Optimisation}} \\
Optimizer        & Adamax \\
Learning rate    & $5{\times}10^{-4}$ \\
Weight decay     & Adamax default \\
Batch size       & 8 (script override; \texttt{train.py} default 12) \\
Epochs           & 100 \\
LR scheduler     & \texttt{ReduceLROnPlateau}
                   (mode=\texttt{max}, factor=0.5, patience=5) \\
Decision metric  & Per-image tolerant mIoU ($\pm$2\,px) \\
Early stopping   & patience = 10 epochs \\
\midrule
\multicolumn{2}{@{}l}{\textit{Loss}} \\
Loss             & \texttt{OGLANetLoss}: six CE terms
                   ($\mathcal{L}_1$--$\mathcal{L}_6$ on $P_1$--$P_6$),
                   equal weights \\
Post-processing  & Remove connected components $<10$\,px \\
\midrule
\multicolumn{2}{@{}l}{\textit{Data}} \\
Input            & $384{\times}384{\times}4$ (RGBC) \\
\midrule
\multicolumn{2}{@{}l}{\textit{Infrastructure}} \\
Hardware         & 1$\times$ NVIDIA A100 40\,GB
                   (NCSA Delta, \texttt{gpuA100x4}) \\
Mixed precision  & Disabled \\
Gradient clip.   & Disabled \\
\bottomrule
\end{tabular}
\end{table}

\subsection{OGLANet + FDA}

\textbf{OGLANet\,+\,FDA.}
FDA is applied identically to MAMNet\,+\,FDA (see \S E.2): per-channel
2-D FFT, central $2b{\times}2b$ amplitude swap with $b{=}\lfloor
\min(H,W)\cdot L\rfloor$, source phase preserved, reconstruction
clipped to $[0,255]$.  Applied as input-level preprocessing in the
dataset pipeline before standard augmentation; model architecture and
optimisation are unchanged.  Run in LOCO mode only; the target pool is
sampled uniformly per iteration from the held-out city's
\texttt{train/images}.

\begin{table}[h]
\centering
\caption{OGLANet\,+\,FDA hyperparameters; settings matching
Table~\ref{tab:oglanet-config} are omitted.}
\label{tab:oglanet-fda}
\small
\begin{tabular}{@{}ll@{}}
\toprule
\textbf{Hyperparameter} & \textbf{Value} \\
\midrule
FDA low-freq.\ ratio $L$ & 0.01 \\
Target pool size         & $\leq 100$ images \\
Target domain            & Held-out city \texttt{train/images} (unlabelled) \\
Target sampling          & Uniform random per iteration \\
Application point        & Input (dataset, pre-augmentation) \\
Modes run                & LOCO only \\
\bottomrule
\end{tabular}
\end{table}

\subsection{OGLANet + SegDeSiC}

\textbf{OGLANet\,+\,SegDeSiC.}
The SegDeSiC module is attached in parallel to \texttt{feat5}
($1024{\times}6{\times}6$, GLAMEncoder output): global average pool
$\to$ Linear(1024, 256) $\to$ ReLU+Dropout(0.3)
$\to$ Linear(256, 256) $\to$ ReLU+Dropout(0.3)
$\to$ Linear(256, 40) $\to$ $\ell_2$ normalisation; weights Xavier
uniform.  GRID encoding ($S{=}10$, $\lambda{\in}[10^{-4},10^{4}]$) and
the cosine-dissimilarity domain loss are identical to
MAMNet\,+\,SegDeSiC (\S E.3).  Total loss combines the deep-supervision
six-CE segmentation loss (equal weights $\tfrac{1}{6}$) with source and
target geographic terms scaled by $\alpha{=}0.5$ (target term
LOCO-only).  Coordinates are not supplied at validation or test time.

\begin{table}[h]
\centering
\caption{OGLANet\,+\,SegDeSiC hyperparameters; settings matching
Table~\ref{tab:oglanet-config} are omitted.}
\label{tab:oglanet-segdesic}
\small
\begin{tabular}{@{}ll@{}}
\toprule
\textbf{Hyperparameter} & \textbf{Value} \\
\midrule
\multicolumn{2}{@{}l}{\textit{SegDeSiC module}} \\
Attachment point         & \texttt{feat5} ($1024{\times}6{\times}6$) \\
GRID scales / range      & 10 / $[10^{-4},10^{4}]$ log-spaced \\
Encoding dim             & 40 \\
MLP hidden dim / dropout & 256 / 0.3 \\
\midrule
\multicolumn{2}{@{}l}{\textit{Loss}} \\
Seg.\ loss               & Mean CE over $P_1$--$P_6$, equal weights \\
Domain loss              & $1-\cos\langle\hat{e},e^{*}\rangle$, batch mean \\
$\alpha$                 & 0.5 (both source and target terms) \\
Target domain (LOCO)     & Held-out city \texttt{train} split (unlabelled) \\
\midrule
\multicolumn{2}{@{}l}{\textit{Optimisation overrides}} \\
Learning rate            & $3{\times}10^{-4}$ (vs.\ $5{\times}10^{-4}$ base) \\
Batch size               & 4 (vs.\ 8 base) \\
Modes run                & LOCO only \\
\bottomrule
\end{tabular}
\end{table}

\subsection{OGLANet + IIM}

\textbf{OGLANet\,+\,IIM.}
The IIM preprocessor (identical to \S E.4: 8 learnable $5{\times}5$
zero-mean kernels, three cross-channel pairs, FuseConv, adaptive II
loss with $r{=}0.01$ and $\gamma\sim U[0.5,2.0]$) is inserted before
the OGLANet stem.  The downstream architecture (GLAMEncoder, DFFM,
decoder, OAM) is unchanged.  The optimiser is switched from Adamax to
Adam to handle the gradient-scale gap between the small IIM module
($\approx$0.008\,M) and the pretrained ResNet-34 backbone
($\approx$21\,M).

\begin{table}[h]
\centering
\caption{OGLANet\,+\,IIM hyperparameters; settings matching
Table~\ref{tab:oglanet-config} are omitted.}
\label{tab:oglanet-iim}
\small
\begin{tabular}{@{}ll@{}}
\toprule
\textbf{Hyperparameter} & \textbf{Value} \\
\midrule
\multicolumn{2}{@{}l}{\textit{IIM module}} \\
Num.\ kernels $K$        & 8 \\
Kernel size              & $5{\times}5$ \\
Cross-channel pairs      & (R,B), (R,G), (G,B) $\to$ 24 features \\
Zero-mean projection     & After every optimizer step \\
\midrule
\multicolumn{2}{@{}l}{\textit{Loss}} \\
II loss                  & SmoothL1 ($\beta{=}1.0$) on IIM features \\
II weighting             & Adaptive, target ratio $r{=}0.01$ \\
Gamma range              & $\gamma\sim U[0.5,2.0]$ \\
\midrule
\multicolumn{2}{@{}l}{\textit{Optimisation overrides}} \\
Optimizer                & Adam (vs.\ Adamax) \\
Learning rate            & $1{\times}10^{-4}$ (vs.\ $5{\times}10^{-4}$) \\
Weight decay             & $1{\times}10^{-4}$ \\
\bottomrule
\end{tabular}
\end{table}

\subsection{OGLANet + ISW}

\textbf{OGLANet\,+\,ISW.}
ISW is applied as in \S E.5 with hooks on \texttt{glam1},
\texttt{glam2}, \texttt{glam3} (the first three GLAMEncoder outputs).
Phase 1 precomputes binary masks
$\widetilde M$ once per LOCO fold from the photometric augmentation
covariance variance, partitioned by $k$-means ($k{=}3$, $m{=}1$);
Phase 2 adds $\lambda\,\mathcal{L}_{\text{ISW}}$ to the six-CE
segmentation loss, with $\mathcal{L}_{\text{ISW}}$ averaged
per-image and across the three layers.  Inference is unchanged.

\begin{table}[h]
\centering
\caption{OGLANet\,+\,ISW hyperparameters; settings matching
Table~\ref{tab:oglanet-config} are omitted.}
\label{tab:oglanet-isw}
\small
\begin{tabular}{@{}ll@{}}
\toprule
\textbf{Hyperparameter} & \textbf{Value} \\
\midrule
$\lambda$ (ISW weight)         & 0.6 \\
Hooked layers                  & \texttt{glam1, glam2, glam3} \\
$k$-means $k$ / $m$            & 3 / 1 \\
Photometric aug.\ (mask precomp.)\!\!
  & \texttt{ColorJitter(0.4,0.4,0.4,0.1)} \\
  & + \texttt{GaussianBlur} $p{=}0.5$ \\
Loss reduction                 & Per-image mean, then layer mean \\
\bottomrule
\end{tabular}
\end{table}

\subsection{OGLANet + MRFP+}

\textbf{OGLANet\,+\,MRFP+.}
HRFP attaches after \texttt{conv1}-\texttt{bn1}-\texttt{relu} of the
GLAMEncoder stem (64\,ch, $H/2$); NP+ attaches at \texttt{glam4}
(512\,ch, $H/32$); HRFP+ adds the max-resolution intermediate HRFP
feature to the decoder's \texttt{s1\_d\_up} stage (64\,ch,
$384{\times}384$).  Modules are toggled independently with $p{=}0.5$.
\emph{The InstanceNorm layers specified in the original MRFP+ paper
are removed} (Option A): the GLAM stages end with two GFEM
self-attention modules whose channel-wise statistics are essential
to subsequent processing, and applying \texttt{InstanceNorm2d}
(affine=False) after GFEM caused class collapse (mIoU $<$1\%) in all
three cities.  HRFP and NP+ alone span both ends of the perturbation
spectrum.  HRFP weights ($\approx$74\,K) are frozen, re-randomised
each forward pass, and discarded at inference.

\begin{table}[h]
\centering
\caption{OGLANet\,+\,MRFP+ hyperparameters; settings matching
Table~\ref{tab:oglanet-config} are omitted.}
\label{tab:oglanet-mrfp}
\small
\begin{tabular}{@{}ll@{}}
\toprule
\textbf{Hyperparameter} & \textbf{Value} \\
\midrule
\multicolumn{2}{@{}l}{\textit{Perturbation modules (training only)}} \\
HRFP attach point        & After stem (64\,ch, $H/2$) \\
NP+ attach point         & \texttt{glam4} (512\,ch, $H/32$) \\
HRFP+ attach point       & Decoder \texttt{s1\_d\_up} (64\,ch, $384{\times}384$) \\
Inst.\ Norm layers       & Removed (Option A) \\
HRFP / NP+ / HRFP+ prob.\ & 0.5 / 0.5 / 0.5 \\
HRFP arch.               & 4 enc + 4 dec, $\sim$1.2$\times$ scale per layer \\
HRFP BN std $\sigma$     & 0.5 \\
\midrule
\multicolumn{2}{@{}l}{\textit{Optimisation overrides}} \\
Learning rate            & $1{\times}10^{-4}$ (vs.\ $5{\times}10^{-4}$) \\
Weight decay             & $1{\times}10^{-4}$ \\
Loss reduction           & Per-image mean CE over $P_1$--$P_6$ \\
\bottomrule
\end{tabular}
\end{table}

\subsection{OGLANet + FADA}

\textbf{OGLANet\,+\,FADA.}
Only the ResNet-34 core inside each GLAM stage
(\texttt{conv1}, \texttt{bn1}, \texttt{layer1}--\texttt{layer4}) is
frozen; the GFEM modules, \texttt{glam5\_conv}, \texttt{glam5\_gfem},
DFFM, decoder, and OAM remain trainable because they are not pretrained
and freezing them collapsed mIoU to $<$1\% in pilot runs.  FADA blocks
are inserted between the frozen LFEM output and the trainable GFEM at
stages 3 and 4 (256\,ch and 512\,ch, $24{\times}24$ and $12{\times}12$);
stages 1, 2, and 5 are excluded for either spatial-cost or
non-pretrained-component reasons.  The block internals (Haar DWT,
LL LoRA-token branch, IN-on-similarity-map HF branch, inverse Haar
DWT) are identical to MAMNet\,+\,FADA (\S E.7).  Trainable parameters:
$\approx$7--9\,M; frozen $\approx$21.3\,M.

\begin{table}[h]
\centering
\caption{OGLANet\,+\,FADA hyperparameters; settings matching
Table~\ref{tab:oglanet-config} are omitted.}
\label{tab:oglanet-fada}
\small
\begin{tabular}{@{}ll@{}}
\toprule
\textbf{Hyperparameter} & \textbf{Value} \\
\midrule
\multicolumn{2}{@{}l}{\textit{FADA configuration}} \\
Adapter stages           & GLAM 3, 4 (LFEM output, pre-GFEM) \\
LoRA rank $r$            & 16 \\
Token length $m$         & 100 \\
Decomposition            & Haar DWT \\
\midrule
\multicolumn{2}{@{}l}{\textit{Optimisation overrides}} \\
Optimizer                & Adam (vs.\ Adamax) \\
Learning rate            & $1{\times}10^{-4}$ (vs.\ $5{\times}10^{-4}$) \\
Weight decay             & $1{\times}10^{-4}$ \\
LR scheduler             & \texttt{ReduceLROnPlateau} (factor=0.5, patience=3) \\
Frozen                   & ResNet-34 core only (BN in eval mode) \\
Trainable                & FADA + GFEM + DFFM + decoder + OAM \\
\bottomrule
\end{tabular}
\end{table}

\subsection{DINOv3}
\label{sec:dinov3-base}

\textbf{Architecture.}
DINOv3 ViT-S/16 backbone (\texttt{dinov3\_vits16}): patch size
$16{\times}16$, embedding dimension $d{=}384$, 12 transformer blocks.
Pretrained weights
\texttt{dinov3\_vits16\_pretrain\_lvd1689m-08c60483.pth}
(\texttt{github.com/facebookresearch/dinov3}); all backbone parameters
trainable (\texttt{frozen\_stages}${=}{-1}$).
Intermediate features extracted after blocks 3, 6, 9, and 11 at
$H/16{\times}W/16$, concatenated along the channel axis, and fused by
a $1{\times}1$-Conv-BN-ReLU ($4{\times}384{\to}384$,
\texttt{bias=False}).  The decoder applies four
\texttt{ConvTranspose2d} ($2{\times}2$, stride 2) stages
($\tfrac{1}{16}{\to}\tfrac{1}{8}{\to}\tfrac{1}{4}{\to}\tfrac{1}{2}{\to}1$)
each followed by Conv-BN-ReLU; channel sequence
$[384, 256, 128, 64, 32]$.  No encoder skip connections (all ViT
features share the same spatial resolution); a final $1{\times}1$ conv
yields 2-class logits.  No auxiliary branches; single output head.
Input is $384{\times}384$ RGB, yielding exactly $24{\times}24$ patches
with no padding; output stride 16.  

\begin{table}[h]
\centering
\caption{DINOv3 training configuration.}
\label{tab:dinov3-config}
\small
\begin{tabular}{@{}ll@{}}
\toprule
\textbf{Hyperparameter} & \textbf{Value} \\
\midrule
\multicolumn{2}{@{}l}{\textit{Optimisation}} \\
Optimizer        & AdamW ($\beta_1{=}0.9$, $\beta_2{=}0.999$) \\
Learning rate    & $5{\times}10^{-5}$
                   (CLI; argparse default $3{\times}10^{-4}$) \\
Weight decay     & 0.05 \\
Batch size       & 8 \\
Epochs (max)     & 100 \\
LR scheduler     & Cosine + linear warmup (5 ep.,
                   $\eta_\text{min}{=}10^{-6}$) \\
Decision metric  & Per-image strict mIoU
                   (\texttt{DetailedEvaluator}) \\
Early stopping   & 20 (single, LOCO); 10 (all) \\
\midrule
\multicolumn{2}{@{}l}{\textit{Loss}} \\
Loss             & \texttt{nn.CrossEntropyLoss} (2-class, no weights,
                   no aux.) \\
Post-processing  & Remove connected components $<10$\,px \\
\midrule
\multicolumn{2}{@{}l}{\textit{Data}} \\
Input size       & $384{\times}384$ \\
Channels         & 3 (RGB) \\
Train aug.       & H-flip $p{=}0.5$, V-flip $p{=}0.5$ (PIL, train split) \\
Normalisation    & ImageNet $\mu, \sigma$ \\
\midrule
\multicolumn{2}{@{}l}{\textit{Infrastructure}} \\
Hardware         & 1$\times$ A100 (NCSA Delta, \texttt{gpuA100x4}) \\
Sampling seed    & 42 (dataset subsampling only) \\
\bottomrule
\end{tabular}
\end{table}

\subsection{DINOv3 + FDA}

\textbf{DINOv3\,+\,FDA.}
FDA is applied as in \S E.2 (per-channel 2-D FFT, central
$2b{\times}2b$ amplitude swap with $b{=}\lfloor\min(H,W)\cdot
L\rfloor$, source phase preserved, reconstruction clipped).
Inserted in the dataset pipeline before PIL$\to$tensor and ImageNet
normalisation; validation and test images are unmodified.  FDA and
SegDeSiC are mutually exclusive in all launch scripts.  Run in LOCO
mode only.

\begin{table}[h]
\centering
\caption{DINOv3\,+\,FDA hyperparameters; settings matching
Table~\ref{tab:dinov3-config} are omitted.}
\label{tab:dinov3-fda}
\small
\begin{tabular}{@{}ll@{}}
\toprule
\textbf{Hyperparameter} & \textbf{Value} \\
\midrule
FDA low-freq.\ ratio $L$ & 0.01 \\
Target pool size         & $\leq 100$ images \\
Target domain            & Held-out city \texttt{train/images} (unlabelled) \\
Application point        & Input (pre-tensor, pre-normalisation) \\
Early-stopping patience  & 10 (all modes) \\
Modes run                & LOCO only \\
\bottomrule
\end{tabular}
\end{table}

\subsection{DINOv3 + SegDeSiC}

\textbf{DINOv3\,+\,SegDeSiC.}
A \texttt{SegDesicModule} is attached to \texttt{feat\_block11}
($384{\times}24{\times}24$, the deepest extracted ViT feature).  The
module predicts a 40-d GRID encoding ($S{=}10$,
$\lambda{\in}[10^{-4},10^{4}]$) via global average pool and a 256-dim
hidden MLP, with cosine-dissimilarity loss against the ground-truth
encoding from $(\varphi,\lambda)$.  In LOCO mode an additional
unlabelled forward pass over the held-out city's training images
contributes a target-domain cosine loss; no segmentation loss is
applied to the target.  Total loss:
$\mathcal{L} = \mathcal{L}_\text{seg} + \alpha(\mathcal{L}^{\text{src}}_\text{dom}
+ \mathcal{L}^{\text{tgt}}_\text{dom})$.

\begin{table}[h]
\centering
\caption{DINOv3\,+\,SegDeSiC hyperparameters; settings matching
Table~\ref{tab:dinov3-config} are omitted.}
\label{tab:dinov3-segdesic}
\small
\begin{tabular}{@{}ll@{}}
\toprule
\textbf{Hyperparameter} & \textbf{Value} \\
\midrule
\multicolumn{2}{@{}l}{\textit{SegDeSiC module}} \\
Attachment point         & \texttt{feat\_block11} ($384{\times}24{\times}24$) \\
GRID scales / range      & 10 / $[10^{-4},10^{4}]$ \\
Encoding dim             & 40 \\
MLP hidden dim / dropout & 256 / 0.3 \\
\midrule
\multicolumn{2}{@{}l}{\textit{Loss}} \\
Domain loss              & $1-\cos\langle\hat{e},e^{*}\rangle$ \\
$\alpha$                 & 0.5 (both terms) \\
Auxiliary weight         & 0 (no aux.\ branches) \\
Target domain (LOCO)     & Held-out city \texttt{train} split (unlabelled) \\
\midrule
\multicolumn{2}{@{}l}{\textit{Optimisation overrides}} \\
Learning rate            & $3{\times}10^{-4}$ (vs.\ $5{\times}10^{-5}$) \\
Batch size               & 4 (vs.\ 8) \\
Early-stopping patience  & 10 (all modes) \\
\bottomrule
\end{tabular}
\end{table}

\subsection{DINOv3 + IIM}

\textbf{DINOv3\,+\,IIM.}
The IIM preprocessor (8 learnable $5{\times}5$ zero-mean kernels,
three cross-channel pairs, FuseConv, adaptive II loss with
$r{=}0.01$, $\gamma\sim U[0.5,2.0]$; identical to \S E.4) is inserted
before the ViT-S/16 stem.  Output is the 3-channel fused image
re-normalised with ImageNet statistics.  The base DINOv3 strict mIoU
decision metric is replaced by tolerant per-image mIoU
($\pm$2\,px) for consistency with the IIM ablation protocol.

\begin{table}[h]
\centering
\caption{DINOv3\,+\,IIM hyperparameters; settings matching
Table~\ref{tab:dinov3-config} are omitted.}
\label{tab:dinov3-iim}
\small
\begin{tabular}{@{}ll@{}}
\toprule
\textbf{Hyperparameter} & \textbf{Value} \\
\midrule
\multicolumn{2}{@{}l}{\textit{IIM module}} \\
Num.\ kernels $K$        & 8 \\
Kernel size              & $5{\times}5$ \\
Cross-channel pairs      & (R,B), (R,G), (G,B) $\to$ 24 features \\
Zero-mean projection     & After every optimizer step \\
\midrule
\multicolumn{2}{@{}l}{\textit{Loss}} \\
II loss                  & SmoothL1 ($\beta{=}1.0$) \\
II weighting             & Adaptive, $r{=}0.01$ \\
Gamma range              & $\gamma\sim U[0.5,2.0]$ \\
\midrule
\multicolumn{2}{@{}l}{\textit{Optimisation overrides}} \\
Decision metric          & Per-image tolerant mIoU ($\pm$2\,px) \\
Early-stopping patience  & 15 epochs \\
\bottomrule
\end{tabular}
\end{table}

\subsection{DINOv3 + ISW}

\textbf{DINOv3\,+\,ISW.}
ISW hooks are registered on the outputs of ViT blocks 3, 6, and 9
(block 11, used by the decoder, is left untouched).  The
\texttt{[B, N+1, D]} token output is stripped of the CLS token and
the $576$ patch tokens reshaped to \texttt{[B, 384, 24, 24]} before
the standardised covariance is computed.  Mask precomputation and
loss formulation are identical to \S E.5.  Inference is unchanged.

\begin{table}[h]
\centering
\caption{DINOv3\,+\,ISW hyperparameters; settings matching
Table~\ref{tab:dinov3-config} are omitted.}
\label{tab:dinov3-isw}
\small
\begin{tabular}{@{}ll@{}}
\toprule
\textbf{Hyperparameter} & \textbf{Value} \\
\midrule
$\lambda$ (ISW weight)         & 0.6 \\
Hooked blocks                  & ViT blocks 3, 6, 9 \\
Token reshape                  & Strip CLS; $576\to24{\times}24$ \\
$k$-means $k$ / $m$            & 3 / 1 \\
Photometric aug.\ (mask precomp.)\!\!
  & \texttt{ColorJitter(0.4,0.4,0.4,0.1)} \\
  & + \texttt{GaussianBlur} $p{=}0.5$ \\
Loss reduction                 & Per-image mean, then block mean \\
Decision metric                & Per-image tolerant mIoU ($\pm$2\,px) \\
Early-stopping patience        & 15 epochs \\
\bottomrule
\end{tabular}
\end{table}

\subsection{DINOv3 + MRFP+}

\textbf{DINOv3\,+\,MRFP+.}
HRFP attaches at \texttt{feat\_block3} ($384{\times}24{\times}24$);
NP+ attaches at \texttt{feat\_block11} ($384{\times}24{\times}24$);
HRFP+ projects the max-resolution intermediate HRFP feature
($\sim$$H/8$, 384\,ch) through a frozen $1{\times}1$ conv
($384{\to}64$) and adds it to the decoder's \texttt{up3} output
(64\,ch, $192{\times}192$).  Modules are toggled independently with
$p{=}0.5$.  As in \S E.13, the original-paper InstanceNorm layers are
removed: in DINOv3 the four block features feed a single
\texttt{Conv}-\texttt{BatchNorm2d}-\texttt{ReLU} fusion, and applying
unconditional IN to blocks 6 and 9 during training corrupted that BN's
running statistics, producing decoder checkerboard artefacts and zero
mIoU at evaluation.  HRFP and NP+ alone span both ends of the
perturbation spectrum.  HRFP+HRFP+ projection adds $\approx$99\,K
frozen parameters; inference is identical to base DINOv3.

\begin{table}[h]
\centering
\caption{DINOv3\,+\,MRFP+ hyperparameters; settings matching
Table~\ref{tab:dinov3-config} are omitted.}
\label{tab:dinov3-mrfp}
\small
\begin{tabular}{@{}ll@{}}
\toprule
\textbf{Hyperparameter} & \textbf{Value} \\
\midrule
\multicolumn{2}{@{}l}{\textit{Perturbation modules (training only)}} \\
HRFP attach point        & \texttt{feat\_block3} (384\,ch, $24{\times}24$) \\
NP+ attach point         & \texttt{feat\_block11} (384\,ch, $24{\times}24$) \\
HRFP+ attach point       & Decoder \texttt{up3} (64\,ch, $192{\times}192$) \\
HRFP+ projection         & Frozen $1{\times}1$ conv $384{\to}64$ \\
Inst.\ Norm layers       & Removed \\
HRFP / NP+ / HRFP+ prob.\ & 0.5 / 0.5 / 0.5 \\
HRFP arch.               & 4 enc + 4 dec, $\sim$1.2$\times$ scale per layer \\
HRFP BN std $\sigma$     & 0.5 \\
\midrule
\multicolumn{2}{@{}l}{\textit{Optimisation overrides}} \\
Decision metric          & Per-image tolerant mIoU ($\pm$2\,px) \\
Early-stopping patience  & 15 epochs \\
\bottomrule
\end{tabular}
\end{table}

\subsection{DINOv3 + FADA}

\textbf{DINOv3\,+\,FADA.}
The DINOv3 ViT-S/16 backbone is fully frozen
(\texttt{requires\_grad=False}); ViT uses LayerNorm so no
\texttt{eval}-mode override is needed.  FADA blocks are inserted after
ViT blocks 3, 6, 9, and 11 (all four feature-extraction blocks), each
operating on the spatial patch tokens \texttt{[B, 384, 24, 24]}: 2-D
Haar DWT, LL LoRA-token branch (content), HF LoRA-token branch with
Instance Normalisation on the similarity map (style suppression),
inverse Haar DWT, and re-injection into the token stream.  Gradients
propagate through frozen blocks (parameters not updated) to upstream
adapters.  Trainable parameters: $\approx$3.4\,M
($4{\times}\!\approx$0.48\,M FADA + $\approx$1.5\,M decoder).

\begin{table}[h]
\centering
\caption{DINOv3\,+\,FADA hyperparameters; settings matching
Table~\ref{tab:dinov3-config} are omitted.}
\label{tab:dinov3-fada}
\small
\begin{tabular}{@{}ll@{}}
\toprule
\textbf{Hyperparameter} & \textbf{Value} \\
\midrule
\multicolumn{2}{@{}l}{\textit{FADA configuration}} \\
Adapter blocks           & ViT blocks 3, 6, 9, 11 \\
LoRA rank $r$            & 16 \\
Token length $m$         & 100 \\
LF branch                & LoRA tokens on LL \\
HF branch                & LoRA tokens + IN on similarity map (LH/HL/HH) \\
Decomposition            & Haar DWT \\
\midrule
\multicolumn{2}{@{}l}{\textit{Optimisation overrides}} \\
Optimizer                & Adam (vs.\ AdamW) \\
Learning rate            & $1{\times}10^{-4}$ (vs.\ $5{\times}10^{-5}$) \\
Weight decay             & $1{\times}10^{-4}$ (vs.\ 0.05) \\
LR scheduler             & \texttt{ReduceLROnPlateau}
                           (factor=0.5, patience=3) \\
Backbone state           & Frozen (LayerNorm; no eval-mode override) \\
Decision metric          & Per-image tolerant mIoU ($\pm$2\,px) \\
Early-stopping patience  & 15 epochs \\
\bottomrule
\end{tabular}
\end{table}

\subsection{CrossEarth}
\label{sec:crossearth-base}

\textbf{CrossEarth.}
We train CrossEarth~\cite{gong2025crossearth} from its official release
on our splits, starting from the release's
DINOv2~\cite{oquab2024dinov2} ViT-L/14 backbone initialisation at patch
size 16.  The backbone is frozen; the trainable components are the Rein
adapters~\cite{wei2024rein} (token length 100, LoRA rank 16) and a
Mask2Former decoder~\cite{cheng2022mask2former} with a two-class head.
The code release omits four components described in the CrossEarth paper
(the Geospatial Semantic Extractor, the cross-attention Injector, the
masked-image-modelling branch with its ASPP decoder, and Earth-Style
Injection), so our cells train the released Rein and Mask2Former
configuration without the Million-AID style pool.  No CrossEarth
per-benchmark checkpoint is used.  Table~\ref{tab:crossearth-config}
lists the configuration and our deviations from the release, which
adapt it to our $384{\times}384$ patches, our $450$-image training
sets, the 4-hour partition limit, and source-only validation.

\begin{table}[h]
\centering
\caption{CrossEarth configuration.  Every CrossEarth cell
(Vanilla, FDA, SIB, BSM; LOCO and within-city upper bound) uses these
settings, high resolution only, seed 0.}
\label{tab:crossearth-config}
\small
\begin{tabular}{@{}ll@{}}
\toprule
\textbf{Hyperparameter} & \textbf{Value} \\
\midrule
\multicolumn{2}{@{}l}{\textit{Architecture}} \\
Backbone                 & DINOv2 ViT-L/14 at patch 16, frozen ($303.7$\,M parameters) \\
Adapters                 & Rein, token length 100, LoRA rank 16 \\
Decoder                  & Mask2Former, 2 classes \\
Trainable parameters     & $23.6$\,M \\
\midrule
\multicolumn{2}{@{}l}{\textit{Optimisation}} \\
Optimizer                & AdamW ($\beta_1{=}0.9$, $\beta_2{=}0.999$, $\epsilon{=}10^{-8}$) \\
Learning rate            & $1{\times}10^{-4}$ \\
Weight decay             & 0.05 \\
LR scheduler             & Polynomial (power 0.9, horizon 40{,}000 iterations) \\
Batch size               & 4 (release: 1) \\
Iterations               & 10{,}000 (release: 30{,}000) \\
Crop                     & $384{\times}384$ (release: $512{\times}512$) \\
Precision                & bf16 autocast; deformable attention in fp32 \\
Validation               & Every 500 iterations, source cities only \\
Checkpoint selection     & Best two-class tolerant mIoU ($\pm$2\,px) on source-city validation \\
Inference                & Whole image (release: sliding window) \\
\midrule
\multicolumn{2}{@{}l}{\textit{Infrastructure}} \\
Hardware                 & 1$\times$ NVIDIA A100 40\,GB (NCSA Delta, \texttt{gpuA100x4}) \\
\bottomrule
\end{tabular}
\end{table}

\textbf{CrossEarth\,+\,FDA.}  FDA is applied as for the other
architectures: low-frequency ratio $L{=}0.01$ ($6{\times}6$ swapped
region per channel), a pool of 100 unlabelled images from the held-out
city's \texttt{train/images}, input stage, training only.

\textbf{CrossEarth\,+\,SIB.}  SIB on CrossEarth is the Haar and VIB path
of Appendix~\ref{app:sib-impl}; FDA preprocessing and skip-attention
gates are CNN-only components and are not applied.  One shared-weight
SIB block is inserted before the Mask2Former pixel decoder and applied
to all four backbone feature scales (strides 4, 8, 16, 32), adding
$2.4$\,M trainable parameters.  The VIB weights, hidden widths ($C/2$
and $C/4$ at $C{=}1024$), KL warm-up and coverage criterion follow
Appendix~\ref{app:sib-impl}; $c^*_{val}$ is fitted on source-city
validation per cell.

\subsection{BSM}
\label{sec:bsm}

\textbf{BSM.}
No official BSM code has been released, so we reimplemented it from
the paper~\cite{luo2023bsm}.  At each training step a frozen AdaIN
style-transfer network~\cite{huang2017adain} re-renders every image in
the batch with the style of another image from the same batch (a random
permutation with distinct pairs), and the result is mixed with the
original at a fixed coefficient $\alpha$.  The style source is the
training batch itself, so no external or held-out-city imagery enters
training.  BSM is applied to the RGB input before the network, with the
architecture's input normalisation inverted before mixing and reapplied
after; the contrast channel of MAMNet and OGLANet passes through
unmodified.  BSM is disabled at evaluation.  Each architecture otherwise
keeps its base configuration (Tables~\ref{tab:mamnet-config},
\ref{tab:oglanet-config}, \ref{tab:dinov3-config},
\ref{tab:crossearth-config}).  The DINOv3/Chicago cell collapsed at
seed 0 (best source-city validation mIoU $5.51$, below the training
floor of $10.0$) and is excluded from Table~\ref{tab:loco-main} without
imputation.

\begin{table}[h]
\centering
\caption{BSM hyperparameters; settings matching each architecture's
base configuration are omitted.}
\label{tab:bsm}
\small
\begin{tabular}{@{}ll@{}}
\toprule
\textbf{Hyperparameter} & \textbf{Value} \\
\midrule
Style network            & AdaIN: VGG-19 encoder to \texttt{relu4\_1} + mirror decoder, frozen \\
Style-network weights    & Public pretrained AdaIN release; 0 trainable parameters \\
Style source             & Permuted training batch \\
Mixing coefficient $\alpha$ & 0.5 (fixed; checked for non-degeneracy on source-city validation) \\
Application probability  & 1.0 per training batch \\
Batch size               & 8 (CrossEarth: 4) \\
Early-stopping patience  & 20 epochs (vs.\ 10; CrossEarth: fixed 10{,}000 iterations) \\
Modes run                & LOCO, high resolution only \\
\bottomrule
\end{tabular}
\end{table}

\subsection{SAR}
\label{sec:sar}

\textbf{SAR.}
SAR~\cite{niu2023sar} trains nothing.  It adapts the normalisation-layer
affine parameters of each cell's Vanilla LOCO checkpoint (for MAMNet,
the matched-450 retrain) on each test image's own prediction, by
minimising the entropy of reliable pixels under sharpness-aware
minimisation~\cite{foret2021sam}.  We run it episodically at batch size
1: parameters are restored before every image, so nothing accumulates
across the target set, and no held-out-city label or source image is
used.  Normalisation layers stay in evaluation mode on source running
statistics, because MAMNet and OGLANet return auxiliary and
deep-supervision outputs in training mode.  We keep upstream's skip list
and collect the affine parameters of every remaining BatchNorm,
LayerNorm and GroupNorm layer.  The two-class entropy constants are
upstream's, rescaled from $C{=}1000$ to $C{=}2$ and never tuned on
target mIoU.  When the reliability filter admits no pixel, the step is
skipped.  Each SAR row in Table~\ref{tab:loco-main} is the Vanilla
cell's per-image values plus the change SAR produces within its own
pass (final step minus step 0), so the row isolates the adaptation
effect from any difference between the SAR pass and the stored Vanilla
inference.

\begin{table}[h]
\centering
\caption{SAR hyperparameters.}
\label{tab:sar}
\small
\begin{tabular}{@{}ll@{}}
\toprule
\textbf{Hyperparameter} & \textbf{Value} \\
\midrule
Adapted parameters       & Affine of BatchNorm/LayerNorm/GroupNorm layers \\
Excluded (upstream)      & \texttt{layer4}, ViT blocks 9--11, final \texttt{norm} \\
Normalisation statistics & Source running statistics (evaluation mode) \\
Optimizer                & SAM ($\rho{=}0.05$) over SGD (momentum 0.9) \\
Learning rate            & $2.5{\times}10^{-4}$ (upstream ViT default, not batch-scaled) \\
Steps per image          & 1 (MAMNet, OGLANet, DINOv3); 10 (CrossEarth) \\
Entropy margin $E_0$     & $0.4\ln 2 = 0.277$ \\
Model-recovery threshold & $0.0201$ nats \\
EMA momentum             & 0.9 \\
Batch size               & 1, episodic \\
\midrule
Hardware                 & 1$\times$ NVIDIA A100 40\,GB (NCSA Delta, \texttt{gpuA100x4}) \\
\bottomrule
\end{tabular}
\end{table}

\textbf{Selection.}  For MAMNet, OGLANet and DINOv3, we selected the
learning rate ($\{1, 10, 100\}\times$ the base rate) and step count
($\{1, 10\}$) on source-city validation in a pilot on the Phoenix fold
that pooled DINOv3 and MAMNet.  Upstream additionally scales the learning
rate by batch size over 64, which at batch size 1 gives
$3.9{\times}10^{-6}$; we do not apply that scaling.

\textbf{SAR on CrossEarth.}  CrossEarth exposes no logits, so SAR
minimises the entropy of the logarithm of the normalised two-class
Mask2Former score.  Its recipe (base rate, 10 steps) is borrowed from
the S-EO pilot, which selected on S-EO source-city validation with
MAMNet, OGLANet and DINOv3 and did not include CrossEarth; no internal
CrossEarth cell selected it.  The SAR pass runs in fp32 while the parent
cells were inferred under bf16 autocast.  Before adaptation, the SAR
pass differs from the stored Vanilla value by $+0.0063$ (Chicago),
$+0.1053$ (Miami) and $+0.0070$ (Phoenix) mIoU; most of the Miami
difference comes from one image with 46 foreground pixels.  This
difference does not enter the row, which adds only the within-pass
change: $-0.0011$, $-0.0054$ and $+0.0020$ mIoU.  The run's divergence
guard, which fires when more than $20\%$ of images predict a foreground
fraction below $0.01$, fires on Miami and Phoenix.  The number of such
images is the same before and after adaptation in every city (11, 42 and
53 of 150), so the flag reflects the parent model.
